Los 15 requisitos no funcionales del backend del SCDEE se enumeran a
continuación. Su detalle ampliado se incluye en \cref{SEC:REQ-NONFUNCTIONAL}.

\begin{nonfunctional}

  \item \textbf{Rendimiento.} \label{rnf:performance}
  Las operaciones \acs{crud} estándar deben completarse en menos de
  500\,ms; las cargas y descargas, en menos de 3\,s; los listados
  paginados, en menos de 1\,s; las exportaciones, en menos de 10\,s; y las
  tareas asíncronas de reconocimiento, en menos de 60\,s por página.

  \item \textbf{Concurrencia.} \label{rnf:concurrency}
  El sistema debe soportar el escenario de 300 alumnos simultáneos en
  consulta tras la publicación de notas, así como al menos 30 correctores
  simultáneos en operación normal de corrección sin degradación apreciable
  del rendimiento.

  \item \textbf{Despliegue.} \label{rnf:deployment}
  El sistema debe poder desplegarse mediante un único fichero
  \texttt{docker-compose.yml} que levante todos los servicios necesarios
  (Django, PostgreSQL, Redis, \textit{workers} de Celery, MinIO,
  \acs{sftp}), con la configuración del entorno externalizada en
  variables de entorno.

  \item \textbf{Escalabilidad.} \label{rnf:scalability}
  La arquitectura de contenedores debe facilitar un futuro despliegue
  distribuido en servicios \textit{cloud} sin cambios arquitectónicos en
  el código fuente, con \textit{workers} replicables horizontalmente y
  aislamiento lógico entre organizaciones por \texttt{organization\_id}.

  \item \textbf{Disponibilidad y observabilidad.}
  \label{rnf:availability}
  El sistema debe exponer un \textit{endpoint} de \textit{health check}
  (\texttt{/health/}) que verifique la conectividad con PostgreSQL, Redis,
  MinIO y los \textit{workers} de Celery, accesible sin autenticación.

  \item \textbf{Seguridad en la \acs{api}.} \label{rnf:api-security}
  El sistema debe implementar \textit{rate limiting} en la \acs{api},
  validar y sanear todas las entradas, e incluir cabeceras \acs{http}
  anti-caché en las respuestas con datos sensibles.

  \item \textbf{Cumplimiento del \acs{rgpd}.} \label{rnf:gdpr}
  El sistema debe cumplir con los requisitos técnicos derivados del
  \acs{rgpd}: cifrado en reposo de datos especialmente sensibles
  (\acs{dni}), auditoría de accesos, \textit{soft-delete} de usuarios y
  almacenamiento seguro de contraseñas.

  \item \textbf{Interoperabilidad.} \label{rnf:interop}
  El \dfn{backend} debe exponer una \acs{api} \acs{rest} versionada bajo
  \texttt{/api/v1/} cuyo contrato se documenta mediante una especificación
  OpenAPI~3.x generada automáticamente, con respuestas en \acs{json} y
  formato uniforme para errores.

  \item \textbf{Mantenibilidad y extensibilidad.}
  \label{rnf:maintainability}
  El sistema debe seguir una arquitectura modular, con autenticador,
  motor de \acs{ocr} y reconocedores implementados como \dfnpl{plugin}
  intercambiables, y un \textit{dispatcher} de reconocedores extensible a
  nuevos tipos de zona sin modificar el código existente.

  \item \textbf{Internacionalización.} \label{rnf:i18n}
  Los contenidos generados para el usuario (correos, notificaciones)
  deben soportar internacionalización; en la versión 1.0, el idioma por
  defecto por organización es el español.

  \item \textbf{Observabilidad y \textit{logging}.}
  \label{rnf:logging}
  El sistema debe mantener dos niveles de \textit{logging} con políticas
  de retención distintas: log de aplicación rotado periódicamente y log
  de auditoría inmutable y permanente.

  \item \textbf{Pruebas.} \label{rnf:testing}
  El sistema debe disponer de pruebas a tres niveles: unitarias para la
  lógica de negocio y utilidades, de integración para la interacción
  entre componentes, y de carga para validar rendimiento y concurrencia.

  \item \textbf{Persistencia e integridad de datos.}
  \label{rnf:persistence}
  La base de datos principal es PostgreSQL y debe garantizar la
  integridad referencial en todo momento; Redis se utiliza como caché,
  \textit{broker} y almacén de la lista negra de \textit{tokens}, no como
  almacenamiento primario.

  \item \textbf{Almacenamiento de objetos.}
  \label{rnf:object-storage}
  MinIO es el servicio de almacenamiento para todos los ficheros
  binarios de tamaño variable (páginas escaneadas, \acs{pdf} en blanco,
  \textit{payloads} de anotaciones), con metadatos en PostgreSQL.

  \item \textbf{Paginación.} \label{rnf:pagination}
  Todos los \textit{endpoints} de listado de la \acs{api} deben devolver
  resultados paginados por \textit{offset}, con tamaño de página
  configurable y metadatos de paginación en cada respuesta.

\end{nonfunctional}
